\begin{dang}{Xác định số vân giao thoa}
$\bullet$ Số vân sáng, vân tối trong vùng giao thoa
\begin{itemize}
\item Số vân sáng (là số lẻ) thoả mãn: \bth{\boder{-\dfrac{L}{2i}\leq k\leq \dfrac{L}{2i}}\ (k\in\mathbb{Z})} \quad hay \quad \boder{N_s=2\sbrace{\dfrac{L}{2i}}+1}
\item Số vân tối (là số chẵn) thoả mãn: \bth{\boder{-\dfrac{L}{2i}-\dfrac{1}{2}\leq k \leq \dfrac{L}{2}-\dfrac{1}{2}}\ (k\in\mathbb{Z})}\quad hay \quad \boder{N_t=2\sbrace{\dfrac{L}{2i}+\dfrac{1}{2}}}
\end{itemize}
$\bullet$ Xác định số vân sáng, vân tối trên MN
\par\textbf{a) Nếu đã có toạ độ $\mathbf{x_M,\ x_N}$ (giả sử $\mathbf{x_M<x_N}$)}
\begin{align*}
\ct{Số vân sáng: }  \boder{\dfrac{x_M}{i}\le k\le \dfrac{x_N}{i}}\ (k\in\mathbb{Z});\quad \ct{Số vân tối: } \boder{\dfrac{x_M}{i}-\dfrac{1}{2}\le k\le \dfrac{x_N}{i}-\dfrac{1}{2}}\ (k\in\mathbb{Z})
\end{align*}
\begin{itemize}
\item M và N cùng phía với vân trung tâm thì \bth{x_M} và \bth{x_N} cùng dấu.
\item M và N khác phía với vân trung tâm thì \bth{x_M} và \bth{x_N} khác dấu.
\item Nếu không tính M và N thì trong các công thức trên thay dấu $"\le"$ bằng dấu $"<"$
\end{itemize}
\par \textbf{b) Nếu biết độ dài $\mathbf{MN = \ell}$}
\begin{bang}[tabularx={X|X|X}]{Số vân sáng và số vân tối}
\[\textbf{Tính chất hai đầu}\]&\[\textbf{Số vân sáng}\]&\[\textbf{Số vân tối}\]\\\hline
\[\textbf{Hai đầu sáng}\]&\[\dfrac{\ell}{i}+1\]&\[\dfrac{\ell}{i}\]\\\hline
\[\textbf{Hai đầu tối}\]&\[\dfrac{\ell}{i}\]&\[\dfrac{\ell}{i}+1\]\\\hline
\[\textbf{Đầu sáng - Đầu tối}\]&\[\dfrac{\ell}{i}+0,5\]&\[\dfrac{\ell}{i}+0,5\]\\\hline
\[\textbf{Đầu sáng - Đầu ?}\]&\[\sbrace{\dfrac{\ell}{i}}+1\]&\[\sbrace{\dfrac{\ell}{i}+0,5}\]\\\hline
\[\textbf{Đầu tối - Đầu ?}\]&\[\sbrace{\dfrac{\ell}{i}+0,5}\]&\[\sbrace{\dfrac{\ell}{i}}+1\]
\end{bang}
\end{dang}
\newpage